\documentclass[11pt]{article}
\usepackage[a4paper,margin=25mm]{geometry}
\usepackage{amsmath,amssymb,amsthm,mathtools,microtype,enumitem}
\usepackage[T1]{fontenc}
\usepackage{lmodern}
\usepackage[colorlinks=true,linkcolor=blue,urlcolor=blue,citecolor=blue]{hyperref}
\newtheorem{theorem}{Theorem}
\newtheorem{proposition}[theorem]{Proposition}
\newtheorem{corollary}[theorem]{Corollary}
\newtheorem{remark}[theorem]{Remark}
\newcommand{\Ric}{\operatorname{Ric}}
\newcommand{\Rm}{\operatorname{Rm}}
\newcommand{\CP}{\mathbb{CP}}
\newcommand{\calC}{\mathcal C}
\title{Q3: Mass Variation and a Cyclic Degeneration Obstruction}
\author{Research continuation based on the supplied October 8 manuscripts}
\date{8 October 2026}
\begin{document}
\maketitle
\begin{abstract}
The fixed objective is a smooth closed, non-Bach-flat global minimizer of
the squared trace-free Ricci functional. We identify a useful freedom in
one surviving degeneration: the marked smooth $\mathcal O(-3)$ filling
admits scalar-flat K\"ahler ALE metrics of either mass sign at the same
curvature integral. Combining this classical existence theory with the
supplied critical-background boundary flux gives a first-order replacement
formula. In a smooth critical orbifold limit, a marked
$\frac13(1,1)$ degeneration of this topology can remain minimizing only
if a specific self-dual curvature eigenvalue vanishes. This is an
obstruction to a degeneration, not a proof of smooth attainment.
\end{abstract}

\section{The fixed existence problem}
For an oriented smooth closed four-manifold, put
\[
 P(g)=\int_M|\Ric_g^\circ|^2\,dV_g,
 \qquad p_*(M)=\inf_{g\in\operatorname{Met}(M)}P(g).
\]
Here Weyl curvature has the Hilbert--Schmidt curvature-operator norm.
The fixed candidate remains $M=X\#\overline{\CP^{2}}$, with $X$ the
classical Godeaux surface, $K_X^2=1$. The preceding work gives
\[
 2\chi(M)+3\tau(M)=0,\qquad
 \frac{16\pi^2}{3}\le p_*(M)<8\pi^2.
 \tag{1}
\]
The lower bound is the classical monopole estimate; the strict upper
bound uses the supplied complete negative-K\"ahler--Einstein insertion
construction. Equation (1) is inherited here.

All conformal classes in this example have negative Yamabe constant.
A smooth $P$-minimizer has constant scalar curvature. If it were
Bach-flat, its remaining Euler--Lagrange equation would give
$s\Ric^\circ=0$, hence it would be Einstein. But $T=2\chi+3\tau=0$
and the Chern--Weil formula would then force $s=W^+=0$, contrary to
the positive monopole scalar budget. Thus \emph{smooth attainment},
without computation of $p_*(M)$, would answer Q3.

The preceding round found an actual embedded $(-3)$-sphere in $X$,
disjoint from the chosen blowup. Its tubular neighbourhood is the
disk bundle $\mathcal O(-3)$, with end $S^3/\mathbb Z_3$. Its usual
scalar-flat K\"ahler ALE metric has $P=8\pi^2/3$; simply comparing this
number with (1) does not exclude a bubble. The present calculation tests
\emph{replacement at finite scale}, rather than adding independent
lower bounds for the core and the bulk.

\section{A fixed filling with arbitrary mass sign}
Let $N$ be the oriented smooth total space of $\mathcal O(-3)$, with
generator $C\in H_2(N;\mathbb Z)$, $C^2=-3$, and scalar end group
$\Gamma=\frac13(1,1)$.

\begin{proposition}[Consequence of classical ALE existence and mass theorems]
There is a Stein complex structure $J$ on this fixed marked smooth
manifold such that, for every $t\in\mathbb R$, an ALE scalar-flat
K\"ahler metric $g_t$ exists with
\[
 \int_C\omega_t=t,\qquad
 c_1(N,J)\cdot C=-1,\qquad
 m(g_t)=-\frac{t}{9\pi},\qquad
 P(g_t)=\frac{8\pi^2}{3}.
 \tag{2}
\]
Mass is normalized as in Hein--LeBrun \cite{HL}.
\end{proposition}

\begin{proof}
In $\mathbb F_1$, write $\Gamma_0^2=-1$, $f^2=0$, $\Gamma_0f=1$.
Choose a smooth section $D\in|\Gamma_0+2f|$. It is ample and has
$D^2=3$. Its complement is affine, hence Stein, and is an affine
$\mathbb C$-bundle with underlying smooth bundle $\mathcal O(-3)$.
Honda's Theorem 1.2 \cite{Honda} realizes this complex structure in
the ALE scalar-flat K\"ahler deformation family of the same smooth
space. The compact generator has image $C=\Gamma_0-f$ in
$H_2(\mathbb F_1)$, so
\[
 C^2=-3,\qquad
 c_1\cdot C=(2\Gamma_0+3f)(\Gamma_0-f)=-1.
\]
Han--Viaclovsky's Theorem 1.13 \cite{HV} explicitly includes the
group $\frac13(1,1)$ and supplies a scalar-flat K\"ahler ALE metric
in \emph{every} K\"ahler class for each complex structure in this
Artin family. Their Remark 6.1 identifies the ALE K\"ahler cone of
this Stein surface with all of $H^2(N;\mathbb R)$. Hence the period
$t$ is unrestricted. Positivity of a K\"ahler form does not require
$t>0$: $C$ is not holomorphic in the Stein complex structure.

The intersection matrix is $Q=[-3]$. Hein--LeBrun's Theorem 4.3 gives
\[
 m=-\frac1{3\pi}(-1)Q^{-1}t=-\frac{t}{9\pi}.
\]
The Ricci form $\varrho$ is anti-self-dual and represents $2\pi c_1$.
The relative square is $c_1^2=-1/3$. Its ALE decay permits the
relative Chern--Weil pairing, and therefore
\[
 P=2\int_N|\varrho|^2=-2\int_N\varrho\wedge\varrho
   =-8\pi^2c_1^2=\frac{8\pi^2}{3}.
\]
The end marking is retained by the smooth bundle identification and
the fixed ALE background used in the all-class existence theorem.
\end{proof}

This proposition imports a concrete freedom from established theory;
it is not a claim that mass-sign flexibility is historically new.
Its relevance is that a replacement metric need not stay within the
radial negative-mass LeBrun family.

\section{The leading tail and its normalization}
Let $J_0$ be the standard complex structure on the end, and set
\[
 B(x)=\frac{|x|^2 I-xx^{\mathsf T}-(J_0x)(J_0x)^{\mathsf T}}{|x|^4}.
 \tag{3}
\]
Han--Viaclovsky, Remark 2.2, supplies differentiated decay
$g-\delta=O(r^{-2})$, $J-J_0=O(r^{-3})$.
Yao's Theorem D(ii) \cite{Yao} then gives a logarithmic K\"ahler
potential whose coefficient is determined by mass, with
$O(r^{-3})$ metric remainder. In suitable end coordinates,
\[
 g_t=\delta+9m(g_t)B+q_t,\qquad
 |\partial^j q_t|\le C_jr^{-3-j}.
 \tag{4}
\]
This assertion concerns a fixed chosen $g_t$; no uniform estimates
over all $t$ are required.

The coefficient in (4) is also checked directly. For $h=aB$,
\[
 \partial_jh_{ij}-\partial_i\operatorname{tr}h
       =\frac{2ax_i}{r^4}.
\]
The ADM flux with coefficient $1/(12\pi^2)$ on
$S^3/\mathbb Z_n$ is $a/(3n)$. For $n=3$, $a=9m$.
Yao's cone-link convention uses $m_Y=3m$ for this end; (4) avoids
any ambiguity from this different convention or from $d^c$.

The $O(r^{-3})$ terms may contribute at later orders, but there is
no independent leading $r^{-2}$ tensor whose sign needs to be
guessed. This is why the mass formula can be used in the full
metric replacement, not merely in a radial model.

\section{Replacement of a smooth critical orbifold point}
Write
\[
 R^+=W^++\frac{s}{12}I,\qquad
 \calC(g)=\int\left(4|W^+|^2+\frac{s^2}{12}\right)dV
          =4\int|R^+|^2dV.
\]
On a closed manifold, $P=\calC-8\pi^2T$, $T=2\chi+3\tau$;
on pieces the corresponding boundary terms must be retained.

\begin{proposition}[First-order replacement; inherited flux normalization]
Let $g$ be a smooth orbifold metric near a point of type
$\frac13(1,1)$, and suppose it is $P$-critical on a neighbourhood
containing the full support of the modification. Let $\zeta$ be the
unit self-dual form singled out by the isotropy representation, and put
\[
 \mu=\langle R^+\zeta,\zeta\rangle_p.
\]
Replace the point by the marked filling $(N,g_t)$ at length scale
$\lambda$. There are smooth matched competitors, equal to $g$
outside that neighbourhood, for which
\[
 P(g_{\lambda,t})=P(g)+\frac{8\pi^2}{3}
       -24\pi^2m(g_t)\mu\lambda^2+o(\lambda^2)
 =P(g)+\frac{8\pi^2}{3}
       +\frac{8\pi}{3}t\mu\lambda^2+o(\lambda^2).
 \tag{5}
\]
Here the linear boundary flux is the corrected critical-background
flux from supplied source A02.59D. The following proof verifies the
transfer to the nonradial filling (4) and the error order.
\end{proposition}

\begin{proof}
In orbifold normal coordinates, $g=\delta+H_2+O(r^3)$, with $H_2$
the quadratic curvature jet. The scaled ALE end is
\[
 g_{t,\lambda}=\delta+9m\lambda^2B+q_{t,\lambda},\qquad
 |\partial^j q_{t,\lambda}|\le C_j\lambda^3r^{-3-j}.
\]
Choose $\rho=\lambda^{3/7}$. On the inner piece use the exact
scalar-flat K\"ahler metric, add the background jet, and extend
that jet smoothly over the core. In unscaled core coordinates this
extension is $O(\lambda^2)$ with bounded derivatives. Its physical
chiral curvature is $O(1)$ on a region of volume $O(\lambda^4)$.
On the end, relative to the common background chiral curvature,
the errors are $O(\lambda^2r^{-2})+O(\lambda^3r^{-3})+O(r)$.

After subtracting the removed bulk integral, the inner error is bounded by
\[
 C\left(\lambda^4+\lambda^2\rho^2+\lambda^3\rho
            +\lambda^4|\log\lambda|+\rho^5\right).
 \tag{6}
\]
This subtraction is essential: the common $O(\rho^4)$ background
term is present in both integrals. Metric norms and volume factors
contribute to the displayed bounds as well.

In the outer region, expand about the exact bulk metric. Its interior
first variation vanishes by criticality. The mass term therefore gives
only the inner boundary flux. The supplied value $-8\pi^2\mu$ for
a unit Burns tail is linear in its coefficient, and quotient integration
divides by three. Thus the leading boundary term is
\[
 -\frac{8\pi^2}{3}(9m)\mu\lambda^2.
\]
Changing the curvature on the inner sphere to its value at $p$ costs
$O(\lambda^2\rho)$. The outer quadratic remainder is
$O(\lambda^4\rho^{-4})$. The differentiated remainder in (4)
has linear flux $O(\lambda^3/\rho)$, squared error
$O(\lambda^6\rho^{-6})$, and mixed error bounded by
$O(\lambda^5\rho^{-5})$. All are $o(\lambda^2)$, as are (6):
for example, $\rho^5=\lambda^{15/7}$ and
$\lambda^4\rho^{-4}=\lambda^{16/7}$.

Finally, the relative $T$ contribution of the filling is $-1/3$.
Consequently the fixed difference between the $P$ and chiral-functional
comparisons is $8\pi^2/3$. This proves (5). The argument uses a fixed
mass and a fixed filling before $\lambda\downarrow0$; it does not
exchange these limits or assume an Einstein bulk.
\end{proof}

\section{The necessary condition for a minimizing degeneration}
\begin{corollary}
Suppose an ideal minimizing configuration has a smooth critical orbifold
bulk with a $\frac13(1,1)$ point, energy equality, and a reconstruction
on the original manifold using a marked $\mathcal O(-3)$ filling.
Then
\[
                    \boxed{\mu=0.}                 \tag{7}
\]
\end{corollary}
\begin{proof}
The relative identity gives $P(N)\ge8\pi^2/3$ for any admissible
filling with this topology and asymptotics, since its chiral energy is
nonnegative. Replacing it by (2) shows that an ideal minimizer must use
exactly this value. If $\mu>0$, use a positive-mass filling; if
$\mu<0$, use a negative-mass filling. For fixed such mass, (5) is
strictly below the ideal energy for all sufficiently small positive
$\lambda$, contradicting minimality on the same smooth manifold.
Other resolved pieces can be recovered afterwards with error smaller
than this fixed improvement.
\end{proof}

The same comparison is local. A full global energy identity can be
replaced by the following precise local hypotheses: the excised regions
in the minimizing sequence have the marked $\mathcal O(-3)$ topology,
their boundary data converge smoothly to the orbifold-ball boundary,
and
\[
 \liminf_jP(U_j)\ge P_g(B_r/\Gamma)+8\pi^2/3.        \tag{8}
\]
For any fixed small $\lambda$, (5) supplies a filling matching $g$
near that boundary with a fixed strict saving. Smooth boundary
convergence permits matching to $g_j$ with error tending to zero.
Leaving the exterior unchanged contradicts minimizing energy. Relative
Chern--Weil and lower semicontinuity of the nonnegative chiral density
yield (8) when the stated marking and boundary convergence are available.
This argument does not require compactness of the exterior.

\section{What is still required for Q3}
For scalar $\mathbb Z_3$ isotropy, the self-dual block at the fixed point is
\[
 R^+=\operatorname{diag}(\mu,\nu,\nu),\qquad
 \mu+2\nu=s/4.
\]
Equation (7) reduces the surviving block to
\[
                  R^+=\operatorname{diag}(0,s/8,s/8).       \tag{9}
\]
For $s<0$ this is a legitimate algebraic curvature block. Negative
scalar curvature alone does not exclude it. In the radial LeBrun
calculation, the attained second-order logarithmic coefficient is
proportional to $-\mu(\mu-s/8)$, so it vanishes in (9). That calculation
must not be assigned without proof to all the nonradial fillings of
Section 2. Their later-order interaction remains to be determined.

The result also does not identify every bubble with $\mathcal O(-3)$,
establish smooth orbifold extension for every selected limit, or exclude
Ricci-flat ALE and open or partially collapsed limits. Those are still
part of the attainment problem. No smooth closed global minimizer has
been constructed in this continuation.

The proved input from the literature is (2)--(4). The new comparison
developed here is (5), using the specified inherited critical flux;
its consequence is (7) under the stated limit and reconstruction
hypotheses. Its concrete use is to replace the former one-sided
restriction $\mu\ge0$ by the equality $\mu=0$, without altering the
core energy or the marked smooth filling.

\begin{thebibliography}{9}
\bibitem{Honda} N. Honda, \emph{Deformation of LeBrun's ALE metrics
with negative mass}, Theorem 1.2.\newline
\url{https://arxiv.org/abs/1204.4857}
\bibitem{HV} J. Han and J. A. Viaclovsky, \emph{Existence and
compactness theory for ALE scalar-flat K\"ahler surfaces}, Forum Math.
Sigma 8 (2020), e1. Theorem 1.13; Remarks 2.2 and 6.1.\newline
\url{https://arxiv.org/abs/1901.05611}
\bibitem{HL} H.-J. Hein and C. LeBrun, \emph{Mass in K\"ahler geometry},
Comm. Math. Phys. 347 (2016), 183--221. Theorem 4.3.\newline
\url{https://arxiv.org/abs/1507.08885}
\bibitem{Yao} Q. Yao, \emph{Mass and Expansion of Asymptotically
Conical K\"ahler Metrics}, Theorem D(ii).\newline
\url{https://arxiv.org/abs/2206.06505}
\bibitem{LeBrun} C. LeBrun, \emph{Curvature Functionals, Optimal
Metrics, and the Differential Topology of 4-Manifolds}, Section 11.\newline
\url{https://arxiv.org/abs/math/0404251}
\bibitem{source} Supplied research manuscripts:
\texttt{01\_Research\_Synthesis.tex} and
\texttt{02\_Mathematical\_Source\_Companion.tex}, 8 October 2026.
In particular, A02.59D for the corrected critical-background flux;
the preceding continuation supplies its general critical-background
form. These are project inputs, distinct from the published sources above.
\end{thebibliography}
\end{document}
